\title{คาบเรียนที่ 6: อนุพันธ์ฟังก์ชันประกอบ และสูตรอนุพันธ์อื่น}
\frame{\titlepage}

\begin{frame}
\frametitle{สารบัญ}
\tableofcontents
\end{frame}


\begin{frame}
\frametitle{ทบทวนอนุพันธ์}
ในคาบเรียนที่ผ่านมาเราได้ศึกษาอนุพันธ์ดังต่อไปนี้
\begin{itemize}
\item $\dx c = 0$ เมื่อ $c$ เป็นค่าคงตัว
\item $\dx x^n = nx^{n-1}$ เมื่อ $n$ เป็นจำนวนจริงใด ๆ 
\item สมบัติอนุพันธ์พีชคณิตของฟังก์ชัน
\begin{eq*}
\dx (f(x) + g(x)) &=& \dx f(x) + \dx g(x) \\
\dx (f(x) - g(x)) &=& \dx f(x) - \dx g(x)\\
\dx (f(x) g(x)) &=& f(x) \dx g(x) + g(x) \dx f(x) \\
\dx \frac{f(x)}{g(x)} &=& \frac{g(x) \dx f(x) - f(x) \dx g(x)}{g(x)^2}
\end{eq*}
\end{itemize}
\end{frame}

\begin{frame}[t]
\begin{ex}
จงหา
\begin{tasks}(3)
\task $\dx 3$
\task $\dx x^4$
\task $\dx (x^2 + 1)$
\end{tasks}
\end{ex}
\begin{sol}
\begin{tasks}(1)
\task $\dx 3 = 0$ เพราะ $3$ เป็นค่าคงตัว
\task $\dx x^4 = 4x^{4-1} = 4x^3$
\task $\dx (x^2+1) = \dx x^2 + \dx 1 = 2x + 0 = 2x$
\end{tasks}
\end{sol}
\end{frame}

%%%%%%%%%%%%%%
\section{สมบัติอนุพันธ์ของฟังก์ชันประกอบ}
\frame{\sectionpage}
%%%%%%%%%%%%%%


\begin{frame}
\frametitle{ทบทวนฟังก์ชันประกอบ}
\begin{thm}{}{}
กำหนดให้ $f(x)$ และ $g(x)$ เป็นฟังก์ชัน แล้วเรานิยามฟังก์ชันประกอบเป็น
\[ (f\circ g)(x) = f(g(x))\]
\end{thm} \pause
\end{frame}

\begin{frame}[t]
\begin{ex}
ให้ $f(x) = x^2$ และ $g(x) = 3x$ จงหา 
\begin{tasks}(2)
\task $(f \circ g)(x)$ 
\task $(g \circ f)(x)$
\end{tasks}
\end{ex}
\begin{sol}
\begin{tasks}(1)
\task
\begin{eqnarray*}
(f \circ g)(x) &=& f(g(x)) \pause \\
&=& (g(x))^2 \pause \\
&=& (3x)^2 \pause = 9x^2
\end{eqnarray*}  \pause
\task
\begin{eqnarray*}
(g \circ f)(x) &=& g(f(x)) \pause \\
&=& 3f(x) \pause \\
&=& 3x^2
\end{eqnarray*}
\end{tasks}
\end{sol}
\end{frame}


\begin{frame}[t]
ในทางปฏิบัติ เราจะเปลี่ยนฟังก์ชันให้เป็นฟังก์ชันประกอบ \pause
\begin{ex}
จงแปลงฟังก์ชัน $f(x) = (x+2)^4$ ให้อยู่ในรูปของฟังก์ชันประกอบ $(u \circ v)(x)$
\end{ex}
\begin{sol}
ให้ $f(x) = (x+2)^4$ \pause 
\begin{itemize}
\item ถ้าเรากำหนดให้ $v(x) = x +2$ \pause จะได้ $f(x) = (v(x))^4$ \pause
\item ถ้าเรากำหนดให้ $u(x) = x^4$ \pause จะได้ 
\[ f(x) = (v(x))^4 \pause = u(v(x)) \pause = (u \circ v)(x)\]
\end{itemize}
\end{sol}
\end{frame}


\begin{frame}[t]
\begin{ex}
จงแปลงฟังก์ชัน $f(x) = (x^2+3x)^3$ ให้อยู่ในรูปของฟังก์ชันประกอบ $(u \circ v)(x)$
\end{ex}
\begin{sol}
ให้ $f(x) = (x^2+3x)^3$ \pause 
\begin{itemize}
\item ถ้าเรากำหนดให้ $v(x) = x^2+3x$ \pause จะได้ $f(x) = (v(x))^3$ \pause 
\item ถ้าเรากำหนดให้ $u(x) = x^3$ \pause จะได้ 
\[ f(x) = (v(x))^3 \pause = u(v(x)) \pause = (u \circ v)(x)\]
\end{itemize}
\end{sol}
\end{frame}


\begin{frame}[t]
\begin{ex}
จงแปลงฟังก์ชัน $f(x) = \frac{1}{(x^3+1)^2}$ ให้อยู่ในรูปของฟังก์ชันประกอบ $(u \circ v)(x)$
\end{ex}
\begin{sol}
ให้ $f(x) = \frac{1}{(x^3+1)^2}$ \pause 
\begin{itemize}
\item ถ้าเรากำหนดให้ $v(x) = x^3 +1$ \pause จะได้ $f(x) = \frac{1}{(v(x))^2}$ \pause
\item ถ้าเรากำหนดให้ $u(x) = \frac{1}{x^2} = x^{-2}$ \pause จะได้ 
\[ f(x) = \frac{1}{(u(x))^2} \pause = (u(x))^{-2} \pause = v(u(x)) \pause= (v \circ u)(x)\]
\end{itemize}
\end{sol}
\end{frame}


\begin{frame}[t]
\begin{ex}
จงเขียนฟังก์ชัน $f(x) = \sqrt[3]{x^5+5}$ ให้อยู่ในรูปของฟังก์ชันประกอบ $(u \circ v)(x)$
\end{ex}
\begin{sol}
ให้ $f(x) = \sqrt[3]{x^5+5}$ \pause 
\begin{itemize}
\item ถ้าเรากำหนดให้ $v(x) = x^5+5$ \pause จะได้ $f(x) = \sqrt[3]{v(x)}$ \pause
\item ถ้าเรากำหนดให้ $u(x) = \sqrt[3]{x} = x^{\frac{1}{3}}$ \pause จะได้ 
\[ f(x) = \sqrt[3]{v(x)} \pause = u(v(x)) \pause = (u \circ v)(x)\]
\end{itemize}
\end{sol}
\end{frame}

\begin{frame}
\frametitle{สูตรอนุพันธ์ของฟังก์ชันประกอบ (กฎลูกโซ่)}\pause 
\begin{thm}{}{}
ถ้า $f(x) = (u \circ v)(x)$ เป็นฟังก์ชันประกอบแล้ว \pause
\begin{eqnarray}
\dx f(x) = \dx (u \circ v)(x) \pause &=& \left(\frac{d}{d v(x)} u(v(x))\right)\left( \dx v(x)\right)
\end{eqnarray}
\end{thm}
โดยที่ความหมายของ $\frac{d}{d v(x)} u(v(x))$ คือ การหาอนุพันธ์ $\frac{d}{dx} u(x)$ แล้วเปลี่ยนตัวแปร $x$ ให้กลายเป็นฟังก์ชัน $v(x)$
\end{frame}


\begin{frame}[t]
\begin{ex}
จงหา $\dx (x+2)^4$
\end{ex}
\begin{sol}
กำหนดให้ $f(x) = (x+2)^4$ \pause เมื่อเขียนให้เป็นรูปของฟังก์ชันประกอบจะได้ \pause
\[ v(x) = x+2\pause, \quad u(x) = x^4 \pause \]
\[ \dx v(x) = 1\pause, \quad \dx u(x) = 4x^3 \rightarrow \frac{d}{d\red{v(x)}} u(\red{v(x)}) = 4(\red{v(x)})^3\pause\]
ฉะนั้น
\begin{eqnarray*}
\dx f(x) = \dx (u \circ v)(x) \pause &=&  \left(\frac{d}{d v(x)} u(v(x))\right)\left( \dx v(x)\right) \pause \\
&=& 4(v(x))^3 (1) \pause \\
&=&  4(x+2)^3
\end{eqnarray*}
\end{sol}
\end{frame}


\begin{frame}[t]
\begin{ex}
จงหา $\dx (x^2+1)^3$
\end{ex}
\begin{sol}
กำหนดให้ $f(x) = (x^2+1)^3$ \pause เมื่อเขียนให้เป็นรูปของฟังก์ชันประกอบจะได้ \pause
\[ v(x) = x^2+1 \pause, \quad u(x) = x^3\pause\]
\[\dx v(x) = 2x\pause, \quad \dx u(x) = 3x^2 \rightarrow \frac{d}{d\red{v(x)}} u(\red{v(x)}) = 3(\red{v(x)})^2\pause\] 
ฉะนั้น
\begin{eqnarray*}
\dx f(x) = \dx (u \circ v)(x) \pause &=&  \left(\frac{d}{d v(x)} u(v(x))\right)\left( \dx v(x)\right) \pause \\
&=& 3(v(x))^2 (2x) \pause \\
&=&  3(x^2+1)^2 (2x) \pause = 6x(x^2+1)^2
\end{eqnarray*}
\end{sol}
\end{frame}


\begin{frame}[t]
\begin{ex}
จงหา $\dx  \frac{1}{(x^3+3x)^4}$
\end{ex}
\begin{sol}
กำหนดให้ $f(x) = \tfrac{1}{(x^3+3x)^4}$ \pause เมื่อเขียนให้เป็นรูปของฟังก์ชันประกอบจะได้ \pause
\[ v(x) = x^3+3x\pause, \quad u(x) = \tfrac{1}{x^4} = x^{-4}\pause\]
\[ \dx v(x) = 3x^2+3\pause, \quad \dx u(x) = -4x^{-5} \rightarrow \frac{d}{dv(x)} u(v(x)) = -4(v(x))^{-5}\pause\]
ฉะนั้น
\begin{eqnarray*}
\dx (u \circ v)(x) \pause &=& \left(\frac{d}{d v(x)} u(v(x))\right)\left( \dx v(x)\right)\pause \\
&=&= -4(v(x))^{-5} \cdot (3x^2+3) \pause \\
&=&  (-4(x^3+3x)^{-5} \cdot (3x^2+3) \pause= -4(3x^2+3)(x^3+3x)^{-5}
\end{eqnarray*}
\end{sol}
\end{frame}

\begin{frame}
\frametitle{สูตรอนุพันธ์ของฟังก์ชันประกอบแบบย่อ}
\begin{thm}{}{}
ถ้า $f(x) = (u \circ v)(x)$ เป็นฟังก์ชันประกอบแล้ว \pause
\begin{eqnarray*}
\dx f(x) = \dx (u \circ v)(x) \pause &=& u'(v(x)) \cdot v'(x) \pause \\
&=& (u'(v) \cdot v' )(x) \nonumber
\end{eqnarray*}
\end{thm}
\end{frame}

\begin{frame}
\frametitle{อนุพันธ์ฟังก์ชันยกกำลัง}
เนื่องจากอนุพันธ์ของฟังก์ชันประกอบในรูปแบบฟังก์ชันยกกำลังใช้บ่อย เราจึงสรุปเป็นสูตรได้ตามนี้
\begin{thm}{}{}
กำหนดฟังก์ชัน $u(x)$ ที่หาอนุพันธ์ได้ และ $n$ เป็นจำนวนจริงที่ไม่เท่ากับ 0 แล้ว
\begin{equation} \label{eq:diff_composite_poly}
\dx (u(x))^n = nu(x)^{n-1} \dx u(x)
\end{equation}
\end{thm}
\end{frame}

\begin{frame}[t]
\begin{ex}
จงหา $\dx (3x+1)^4$
\end{ex}
\begin{sol}
\begin{eq*}
\dx (3x+1)^4 &=& 4(3x+1)^3 \dx (3x+1) \\
&=& 4 (3x+1)^3 (3) \\
&=& 12 (3x+1)^3
\end{eq*}
\end{sol}
\end{frame}


\begin{frame}[t]
\begin{ex}
จงหา $\dx (x^2-1)^5$
\end{ex}
\begin{sol}
\begin{eq*}
\dx (x^2-1)^5 &=& 5(x^2-1)^4 \dx (x^2-1) \\
&=& 5 (x^2-1)^4 (2x) \\
&=& 10x (x^2-1)^4
\end{eq*}
\end{sol}
\end{frame}


\begin{frame}[t]
\begin{ex}
จงหา $\dx \sqrt{6x+5}$
\end{ex}
\begin{sol}
\begin{eq*}
\dx \sqrt{6x+5} &=& \dx (6x+5)^{1/2} \\
&=& \frac{1}{2} (6x+5)^{-1/2} \dx (6x+5) \\
&=& \frac{1}{2 \sqrt{ 6x+5}} (6) \\
&=& \frac{3}{\sqrt{6x+5}}
\end{eq*}
\end{sol}
\end{frame}



%%%%%%%%%%%%%%
\section[ฟังก์ชันเลขชี้กำลัง]{อนุพันธ์ของฟังก์ชันเลขชี้กำลัง}
\frame{\sectionpage}
%%%%%%%%%%%%%%

\begin{frame}
\frametitle{ฟังก์ชันเลขชี้กำลัง}
ฟังก์ชันนึงที่มีความสำคัญมากในการสร้างแบบจำลองทางคณิตศาสตร์ คือฟังก์ชันเลขชี้กำลัง
\[ f(x) = a^x\]
โดยที่ $a$ เป็นจำนวนจริงบวก

\begin{remark}{ข้อควรระวัง}
ฟังก์ชันยกกำลัง $f(x) = x^n$ และฟังก์ชันเลขชี้กำลัง $f(x) = n^x$ เป็นฟังก์ชันที่แตกต่างกัน โดยสังเกตได้จากตำแหน่งของตัวแปร เช่น
\begin{eq*}
f(x) = x^4 &\then& f(3) = 3^4 = (3)(3)(3)(3) = 81 \\
g(x) = 4^x &\then& g(3) = 4^3 = (4)(4)(4) = 64
\end{eq*}
\end{remark}
\end{frame}

\begin{frame}
\frametitle{ฟังก์ชันลอการิทึม}
ฟังก์ชันที่ต้องพิจารณาคู่กับฟังก์ชันเลขชี้กำลังเสมอ คือฟังก์ชันลอการิทึม ซึ่งเป็นฟังก์ชันผกผันของฟังก์ชันเลขชี้กำลัง และมีนิยามดังนี้
\begin{defn}{}{}
นิยามฟังก์ชันลอการิทึมฐานธรรมชาติ 
\[f(x) = \ln x\]
เป็นฟังก์ชันผกผันของ $e^x$ \pa นั่นคือ
\[ y = \ln x \text{ ก็ต่อเมื่อ } e^y = x\]
\end{defn}
เช่น ถ้ากำหนดให้ $a = \ln (e^7)$ แล้วจะได้ว่า $e^a = e^7$ นั่นคือ $a=7$ นั่นเอง
\end{frame}

\begin{frame}
\frametitle{อนุพันธ์ฟังก์ชันชี้กำลัง}
\begin{thm}{}{}
นิยามฟังก์ชัน $f(x) = e^x$ \pa เมื่อ $e$ เป็นค่าคงตัว $e \approx 2.71827$ \pa
แล้วจะได้ว่า
\begin{equation} \label{eq:diff_exp}
\dx e^x = e^x\
\end{equation}
\pa กำหนดให้ $u = u(x)$ เป็นฟังก์ชันที่หาอนุพันธ์ได้ \pa จากกฎลูกโซ่ของฟังก์ชันประกอบจะได้ว่า\pa 
\[ \dx e^{u(x)} = e^{u(x)} \dx u(x)\]
\end{thm}
\end{frame}


\begin{frame}[t]
\begin{ex}
จงหา
\begin{tasks}(3)
\task $\dx e^{2x}$
\task $\dx e^{3x+4}$
\task $\dx e^{x^2}$
\end{tasks}
\end{ex}
\begin{sol}
\begin{eqnarray*}
\dx e^{2x} \pa = e^{2x} \dx 2x \pa &=& e^{2x} (2) \pa \\
&=& 2e^{2x} \pa \\
\dx e^{3x+4} \pa = e^{3x+4} \dx (3x+4) \pa &=& e^{3x+4} (3) \pa \\
&=& 3e^{3x+4} \pa \\
\dx e^{x^2} = e^{x^2} \dx x^2 \pa &=& e^{x^2} (2x) \pa \\
&=& 2xe^{x^2}
\end{eqnarray*}
\end{sol}
\end{frame}


\begin{frame}[t]
\begin{ex}
จงหา
\begin{tasks}(2)
\task $\dx xe^x$
\task $\dx e^{e^x}$
\end{tasks}
\end{ex}
\begin{sol}
\begin{eqnarray*}
\dx xe^x \pa &=& x \dx e^x + e^x \dx x \mathnote{อนุพันธ์ของผลคูณฟังก์ชัน} \pa \\
&=& x(e^x) + e^x(1) \pa \\
&=& xe^x + e^x \pa \\
\dx e^{e^x} \pa &=& e^{e^x} \dx e^x \pa \\
&=& e^{e^x} (e^{x}) \pa \\
&=& e^{(e^x + x)} \mathnote{สมบัติของฟังก์ชันเลขชี้กำลัง}
\end{eqnarray*}
\end{sol}
\end{frame}

\begin{frame}
\frametitle{ฟังก์ชันเลขชี้กำลัง สำหรับฐานใด ๆ}
เราศึกษาอนุพันธ์ของฟังก์ชัน $f(x) = e^x$ เนื่องจากมีสมบัติที่จำได้ง่ายว่า
\[ \dx e^x = e^x\]
ถ้าฟังก์ชันเลขชี้กำลังมีฐานอื่นที่ไม่ใช่ $e$ เราสามารถใช้สมบัติของฟังก์ชันเลขชี้กำลังเพื่อแปลงได้ดังนี้ กำหนดฟังก์ชัน $f(x) = a^x$ แล้วจะได้ว่า
\[a^x = (e^{(\ln a)})^x = e^{ (\ln a) x} \]
โดยที่ $\ln a$ เป็นค่าคงที่ ถ้าเรามองเลขชี้กำลังเป็นฟังก์ชัน $u(x) = (\ln a)x$ และใช้กฎลูกโซ่จะได้
\begin{eq*}
\dx a^x \pa = \dx \left(e^{(\ln a)x}\right) \pa &=& e^{(\ln a)x} \left( \dx (\ln a) x\right) \\
\pa &=& e^{(\ln a)x} (\ln a) \pa = (\ln a) a^x
\end{eq*}
\end{frame}

\begin{frame}
\frametitle{อนุพันธ์ฟังก์ชันเลขชี้กำลัง สำหรับฐานใด ๆ}
\begin{thm}{}{}
กำหนดฟังก์ชัน $f(x) = a^x$ แล้วจะได้ว่า
\begin{equation} \label{eq:diff_exp_base_a}
\dx a^x = (\ln a)a^x
\end{equation}
เมื่อ $\ln a$ เป็นค่าคงตัว
\pa กำหนดให้ $u = u(x)$ เป็นฟังก์ชันที่หาอนุพันธ์ได้ \pa จากกฎลูกโซ่ของฟังก์ชันประกอบจะได้ว่า\pa 
\[ \dx a^{u(x)} = (\ln a)a^{u(x)} \dx u(x)\]
\end{thm}
\end{frame}


\begin{frame}[t]
\begin{ex}
จงหา
\begin{tasks}(3)
\task $\dx 2^{x}$
\task $\dx \left(\frac{1}{3^x}\right)$
\task $\dx 4^{x^2}$
\end{tasks}
\end{ex}
\begin{sol}
\begin{eqnarray*}
\dx 2^x \pa &=& (\ln 2)2^x \\
\dx \left(\frac{1}{3^x}\right) &=& \dx (3^{-x}) \pa = (\ln 3)3^{-x} \left( \dx (-x) \right) \\ 
\pa &=& -(\ln 3)3^{-x} \\
\dx 4^{x^2} &=& (\ln 4) 4^{x^2} \left( \dx x^2\right) \\
\pa &=& (2 \ln 4)(x)4^{x^2}
\end{eqnarray*}
\end{sol}
\end{frame}

%%%%%%%%%%%%%%
\section[ลอการิทึม]{อนุพันธ์ของฟังก์ชันลอการิทึม}
\frame{\sectionpage}
%%%%%%%%%%%%%%

\begin{frame}
\frametitle{อนุพันธ์ฟังก์ชันลอการิทึม}
\begin{thm}{}{}
กำหนดให้ $f(x) = \ln x$ แล้วจะได้ว่า \pa 
\begin{equation} \label{eq:diff_ln}
\dx \ln x = \frac{1}{x}
\end{equation}
\pa กำหนดให้ $u = u(x)$ เป็นฟังก์ชันที่หาอนุพันธ์ได้ จากกฎลูกโซ่ของฟังก์ชันประกอบจะได้ว่า \pa
\[ \dx \ln (u(x)) = \frac{1}{u(x)} \left(  \dx u(x) \right)\]
\end{thm}
\end{frame}



\begin{frame}[t]
\begin{ex}
จงหา
\begin{tasks}(3)
\task $\dx \ln (3x)$
\task $\dx \ln (4x-2)$
\task $\dx \ln(x^3+2x)$
\end{tasks}
\end{ex}
\begin{sol}
\begin{eqnarray*}
\dx \ln (3x) \pa &=& \frac{1}{3x} \left(\dx 3x\right) \pa = \left(\frac{1}{3x} \right)( 3 )= \frac{1}{x} \pa \\
\dx \ln (4x-2) \pa &=& \frac{1}{4x-2} \left(\dx (4x-2)\right) \pa = \left(\frac{1}{4x-2}\right) (4) = \frac{4}{4x-2} \pa \\
\dx \ln (x^3+2x) \pa &=& \frac{1}{x^3+2x} \left(\dx (x^3+2x)\right) \pa = \left(\frac{1}{x^3+2x}\right) (3x^2+2) \pa \\
&=& \frac{3x^2+2}{x^3+2x}
\end{eqnarray*}
\end{sol}
\end{frame}


\begin{frame}[t]
\begin{ex}
จงหา $\dx \ln(e^x + x^e)$
\end{ex}
\begin{sol}
\begin{eqnarray*}
\dx \ln (e^x + x^e) \pa &=& \frac{1}{e^x+x^e} \left( \dx (e^x + x^e) \right) \pa \\
&=& \left( \frac{1}{e^x+x^e} \right) (e^x + ex^{e-1}) \pa \\
&=& \frac{e^x + ex^{e-1}}{e^x + x^e} \pa \\
\end{eqnarray*}
\end{sol}
\end{frame}


\begin{frame}[t]
\begin{ex}
จงหา $\dx \ln (\ln x)$
\end{ex}
\begin{sol}
\begin{eqnarray*}
\dx \ln (\ln x) \pa &=& \frac{1}{\ln x} \left( \dx \ln x\right) \pa \\
&=& \frac{1}{\ln x} \left( \frac{1}{x}\right) \pa \\
&=& \frac{1}{x\ln x}
\end{eqnarray*}
\end{sol}
\end{frame}

\begin{frame}
\frametitle{ข้อควรระวัง 1}
\begin{remark}{}
การเขียนฟังก์ชันลอการิทึมต้องระวังเรื่องการใส่วงเล็บ พิจารณา
\[ f(x) = \ln x + 1\]
ฟังก์ชันนี้สามารถตีความได้สองอย่างคือ $(\ln x) + 1$ หรือ $\ln (x+1)$ ดังนั้น เพื่อหลีกเลี่ยงความกำกวม จึงควรจะใส่วงเล็บทุกครั้ง\textbf{ถ้า}ตัวแปรในฟังก์ชันลอการิทึมไม่ใช่ $x$ โดด ๆ เช่น
\[ \ln (2x), \quad \ln(x^4), \quad \ln (x^2+1), \quad \ln\left( \tfrac{1}{x} \right) \]
และถ้ามีการบวกลบหรือคูณฟังก์ชันอื่น ควรจะวางไว้หน้าฟังก์ชัน $\ln$ เช่น
\[ 2 \ln x, \quad x^3 + \ln x, \quad x\ln (x^2)\]
\end{remark}
\end{frame}

\begin{frame}
\frametitle{ข้อควรระวัง 2}
\begin{remark}{}
ฟังก์ชัน $f(x) =(\ln x)^2$ และ $g(x) = \ln (x^2)$ เป็นคนละฟังก์ชันกัน
\begin{itemize}
\item $(\ln x)^2$ หมายความว่าให้หาค่าของ $\ln x$ ก่อน จากนั้นนำค่าที่ได้ไปยกกำลังสอง
\item $\ln (x^2)$ หมายความว่าให้นำ $x$ ไปยกกำลังสองก่อน จากนั้นนำค่าที่ได้ไปใส่ในฟังก์ชันลอการิทึม
\end{itemize}
ฉะนั้น เวลาใช้กฎลูกโซ่จึงต้องเรียงลำดับให้ถูกต้อง ดังตัวอย่างต่อไปนี้

\end{remark}
\end{frame}


\begin{frame}[t]
\begin{ex}
จงหา 
\begin{tasks}(2)
\task $\dx (\ln x)^2$
\task $\dx \ln (x^2)$
\end{tasks}
\end{ex}

\begin{sol} 
\begin{columns}
\begin{column}{0.5\textwidth}
กำหนดให้ $u = \ln x$ จะได้ \pa
\begin{eq*}
\dx (\ln x)^2  \pa &=& \dx (u^2) \\
\pa &=& 2u \dx (u) \\
\pa &=& 2\ln x \left( \dx \ln x \right) \\
\pa &=& 2 \ln x \left( \frac{1}{x} \right) \pa = \frac{2\ln x}{x}
\end{eq*}
\end{column}
\begin{column}{0.5\textwidth}
กำหนดให้ $u = x^2$ จะได้ \pa
\begin{eq*}
\dx \ln (x^2)  \pa &=& \dx \ln u \\
\pa &=& \frac{1}{u} \dx (u) \\
\pa &=& \frac{1}{x^2} \left( \dx x^2\right) \\
\pa &=& \frac{1}{x^2} (2x) \pa = \frac{2x}{x^2} \pa = \frac{2}{x}
\end{eq*}
\end{column}
\end{columns}
\end{sol}

\end{frame}

\begin{frame}
\frametitle{ฟังก์ชันลอการิทึม สำหรับฐานใด ๆ}
เราได้ศึกษาอนุพันธ์ฟังก์ชันลอการิทึมฐานธรรมชาติ
\[ \dx \ln x = \frac{1}{x}\]
\pa ซึ่งที่จริงแล้วเราสามารถหาอนุพันธ์ของฟังก์ชันลอการิทึมฐานใดก็ได้ \pa โดยใช้สมบัติของฟังก์ชันลอการิทึมที่ว่า
\[\log_a x \pa = \frac{\ln x}{\ln a}\]
\pa ฉะนั้น 
\begin{eq*}
\dx (\log_a) x \pa = \dx \left(\frac{\ln x}{\ln a}\right) \pa &=& \frac{1}{\ln a} \left( \dx \ln x \right) \mathnote{$\ln a$ เป็นค่าคงตัว} \\
\pa &=& \frac{1}{\ln a} \left( \frac{1}{x} \right) \pa = \frac{1}{(\ln a)x}
\end{eq*}
\end{frame}

\begin{frame}
\frametitle{อนุพันธ์ฟังก์ชันลอการิทึม สำหรับฐานใด ๆ}
\begin{thm}{}{}
กำหนดให้ $f(x) = \log_a x$ แล้วจะได้ว่า \pa 
\begin{equation} \label{eq:diff_ln}
\dx (\log_a) x = \frac{1}{(\ln a)x}
\end{equation}
\pa กำหนดให้ $u = u(x)$ เป็นฟังก์ชันที่หาอนุพันธ์ได้ จากกฎลูกโซ่ของฟังก์ชันประกอบจะได้ว่า \pa
\[ \dx \log_a (u(x)) = \frac{1}{(\ln a)u(x)} \left(\dx u(x)\right)\]
\end{thm}
\end{frame}


%%%%%%%%%%%%%%
\section[ตรีโกณมิติ]{อนุพันธ์ของฟังก์ชันตรีโกณมิติ}
\frame{\sectionpage}
%%%%%%%%%%%%%%

\begin{frame}
\frametitle{อนุพันธ์ฟังก์ชัน $f(x) = \sin x$}
\begin{thm}{}{}
กำหนดให้ $f(x) = \sin x$ \pa จะได้ว่า
\begin{equation} 
\dx \sin x = \cos x
\end{equation}
\pa กำหนดให้ $u = u(x)$ เป็นฟังก์ชันที่หาอนุพันธ์ได้ จากกฎลูกโซ่ของฟังก์ชันประกอบจะได้ว่า \pa
\[ \dx \sin (u(x)) = \cos u(x) \left( \dx u(x)\right)\]
\end{thm}
\end{frame}


\begin{frame}[t]
\begin{ex}
จงหา
\begin{tasks}(3)
\task $\dx \sin (4x)$
\task $\dx \sin (5x+1)$ 
\task $\dx \sin (e^x)$
\end{tasks}
\end{ex}
\begin{sol}
\begin{eqnarray*}
\dx \sin (4x) \pa &=& \cos (4x) \dx 4x \pa \\
&=& (\cos (4x))(4) \pa = 4 \cos (4x) \pa \\
\dx \sin (5x+1) \pa &=& \cos (5x+1) \dx (5x+1) \pa \\
&=& 5 \cos (5x+1) \pa \\
\dx \sin (e^x) \pa &=& \cos (e^x) \dx (e^x) \pa \\
&=& e^x \cos (e^x)
\end{eqnarray*}
\end{sol}
\end{frame}

\begin{frame}
\frametitle{อนุพันธ์ฟังก์ชัน $f(x) = \cos x$}
\begin{thm}{}{}
กำหนดให้ $f(x) = \cos x$ \pa จะได้ว่า
\begin{equation}
\dx \cos x = -\sin x
\end{equation} 
\pa กำหนดให้ $u = u(x)$ เป็นฟังก์ชันที่หาอนุพันธ์ได้ จากกฎลูกโซ่ของฟังก์ชันประกอบจะได้ว่า
\[ \dx \cos (u(x)) = -\sin u(x) \left( \dx u(x) \right)\]
\end{thm}
\end{frame}


\begin{frame}[t]
\begin{ex}
จงหา
\begin{tasks}(3)
\task $\dx \cos (-2x)$
\task $\dx \cos (x^2+2)$ 
\task $\dx \cos (\ln x)$
\end{tasks}
\end{ex}
\begin{sol}
\begin{eqnarray*}
\dx \cos (-2x) \pa &=& -\sin (-2x) \left(\dx (-2x)\right)\pa \\
&=& (-\sin(-2x))  (-2) \pa = 2 \sin (-2x) \pa \\
\dx \cos (x^2+2) \pa &=& -\sin (x^2+2) \left(\dx (x^2+2) \right) \pa \\
&=& (-\sin(x^2+2))(2x) \pa = -2x\sin(x^2+2) \pa \\
\dx \cos (\ln x) \pa &=& -\sin (\ln x) \left(\dx \ln x\right) \pa = (-\sin (\ln x))\left( \frac{1}{x} \right) \\
\pa &=& \tfrac{-\sin (\ln x)}{x}
\end{eqnarray*}
\end{sol}
\end{frame}

\begin{frame}
\frametitle{อนุพันธ์ฟังก์ชัน $f(x) = \tan x$}
\begin{thm}{}{}
กำหนดให้ $f(x) = \tan x$ จะได้ว่า
\begin{equation}
\dx \tan x = (\sec x)^2
\end{equation}
กำหนดให้ $u = u(x)$ เป็นฟังก์ชันที่หาอนุพันธ์ได้ จากกฎลูกโซ่ของฟังก์ชันประกอบจะได้ว่า
\[ \dx \tan (u(x)) = ( \sec u(x) )^2 \left( \dx u(x) \right)\]
\end{thm}
\end{frame}

\begin{frame}
\frametitle{อนุพันธ์ฟังก์ชันตรีโกณมิติอื่น}
สำหรับฟังก์ชันตรีโกณมิติส่วนกลับ ได้แก่
\[ \sec x = \frac{1}{\cos x}, \qquad \csc x = \frac{1}{\sin x}, \qquad \cot x = \frac{1}{\tan x}\]
เราสามารถหาสูตรได้โดยใช้อนุพันธ์ฟังก์ชันผลหาร ซึ่งจะได้ผลลัพธ์ดังนี้
\begin{thm}{}{}
\begin{eq}
\dx \sec u(x) &=& \sec u(x) \tan u(x) \dx u(x) \\
\dx \csc u(x) &=& -\csc u(x) \cot u(x) \dx u(x) \\
\dx \cot u(x) &=& -(\csc u(x))^2 \dx u(x)
\end{eq}
\end{thm}
\end{frame}


\begin{frame}[t]
\begin{ex}
จงหา
\begin{tasks}(3)
\task $\dx \tan (3x+2)$
\task $\dx \sec (x^3)$ 
\task $\dx \cot (\cos x)$
\end{tasks}
\end{ex}
\begin{sol}
\begin{eqnarray*}
\dx \tan(3x+2) \pa &=& (\sec (3x+2))^2 \left( \dx (3x+2) \right) = 3(\sec (3x+2)^2) \\
\dx \sec (x^3) \pa &=& \sec (x^3) \tan (x^3) \left( \dx (x^3) \right) \\
\pa &=& 3x^2 \sec (x^3) \tan (x^3) \\
\dx \cot (\cos x) \pa &=& -(\csc (\cos x))^2 \left( \dx \cos x\right) \\
\pa &=& - (\csc (\sin x))^2 (- \sin x) \pa = (\csc (\cos x))^2 (\sin x)
\end{eqnarray*}
\end{sol}
\end{frame}

%%%%%%%%%%%%%%
\section[ตรีโกณมิติผกผัน]{อนุพันธ์ของฟังก์ชันตรีโกณมิติผกผัน}
\frame{\sectionpage}
%%%%%%%%%%%%%%

\begin{frame}
\frametitle{ฟังก์ชันตรีโกณมิติผกผัน}
ถ้ากำหนดโดเมนที่เหมาะสม แล้วฟังก์ชันตรีโกณมิติทั้ง 6 ฟังก์ชัน ได้แก่
\[\sin x, \cos x, \tan x, \csc x, \sec x, \cot x\]
จะมีฟังก์ชันผกผัน ซึ่งเรานิยามด้วย
\[\arcsin x, \arccos x, \arctan x, \arccsc x, \arcsec x, \arccot x\]
ตามลำดับ ในบางตำราจะเขียนเป็น
\[\sin^{-1} x, \cos^{-1} x, \tan^{-1} x, \csc^{-1} x, \sec^{-1} x, \cot^{-1} x\]
แต่เพื่อไม่ให้สับสนเรื่องฟังก์ชันผกผัน $\sin^{-1} x$ กับฟังก์ชันตรีโกณมิติยกกำลัง $\sin^{n} x$ เราจึงหลีกเหลี่ยงการใช้สัญลักษณ์ $^{-1}$ และตั้งชื่อฟังก์ชันผกผันด้วยการใส่คำว่า arc นำหน้าแทน
\end{frame}

\begin{frame}
\frametitle{อนุพันธ์ฟังก์ชันตรีโกณมิติผกผัน}
เราหาอนุพันธ์ของฟังก์ชันตรีโกณมิติผกผันด้วยเทคนิคเดียวกัน โดยจะได้ผลลัพธ์ดังนี้
\begin{thm}{}{}
\begin{columns}
	\begin{column}{0.5\textwidth}
			\begin{itemize}
				\item $\dx \arcsin x = \frac{1}{\sqrt{1-x^2}}$
				\item $\dx \arccos x = -\frac{1}{\sqrt{1-x^2}}$
				\item $ \dx \arctan x = \frac{1}{\sqrt{1+x^2}}$		
			\end{itemize}
	\end{column}
	\begin{column}{0.5\textwidth}
			\begin{itemize}
				\item $\dx \arccot x = -\frac{1}{\sqrt{1+x^2}}$
				\item $\dx \arcsec x = \frac{1}{|x|\sqrt{1-x^2}}$
				\item $ \dx \arccsc x = -\frac{1}{|x|\sqrt{1-x^2}}$
			\end{itemize}
	\end{column}
\end{columns}
\end{thm}
บทพิสูจน์ของทฤษฎีบทนี้ อยู่ในส่วนบทพิสูจน์ F)
\end{frame}

\begin{frame}
\frametitle{อนุพันธ์ฟังก์ชันตรีโกณมิติผกผัน (ต่อ)}
กำหนดให้ $u = u(x)$ เป็นฟังก์ชันที่หาอนุพันธ์ได้ จะได้สูตรของฟังก์ชันประกอบดังนี้
\begin{thm}{}{}
\begin{columns}
	\begin{column}{0.5\textwidth}
			\begin{itemize}
				\item $\dx \arcsin u = \frac{1}{\sqrt{1-u^2}} \dux $
				\item $\dx \arccos u = -\frac{1}{\sqrt{1-x^2}} \dux $
				\item $ \dx \arctan u = \frac{1}{\sqrt{1+x^2}} \dux$		
			\end{itemize}
	\end{column}
	\begin{column}{0.5\textwidth}
			\begin{itemize}
				\item $\dx \arccot u = -\frac{1}{\sqrt{1+x^2}} \dux $
				\item $\dx \arcsec u = \frac{1}{|x|\sqrt{1-x^2}} \dux $
				\item $ \dx \arccsc u = \frac{-1}{|x|\sqrt{1-x^2}} \dux$
			\end{itemize}
	\end{column}
\end{columns}
\end{thm}
\end{frame}


\begin{frame}[t]
\begin{ex}
จงหา
\begin{tasks}(3)
\task $\dx \arcsin (3x)$
\task $\dx \arccos (x^4)$ 
\task $\dx \arctan (\sqrt{x})$
\end{tasks}
\end{ex}
\begin{sol}
\begin{eqnarray*}
\dx \arcsin (\red{3x}) \pa &=& \frac{1}{\sqrt{1-(\red{3x})^2}} \dx (\red{3x}) \pa = \frac{3}{\sqrt{1-9x^2}}\\
\dx \arccos (\red{x^2}) \pa &=& -\frac{1}{\sqrt{1-(\red{x^2})^2}}\dx (\red{x^2}) \pa = -\frac{2x}{\sqrt{1-x^4}} \pa \\
\dx \arctan (\red{\sqrt{x}}) \pa &=& \frac{1}{\sqrt{1 + (\red{\sqrt{x}})^2}} \dx (\red{\sqrt{x}}) \pa \\
&=& \frac{1}{\sqrt{1+x}} \left(\frac{1}{2}x^{-\frac{1}{2}} \right) = \frac{1}{2\sqrt{x}\sqrt{1+x}} = \frac{1}{2\sqrt{x(1+x)}}
\end{eqnarray*}
\end{sol}
\end{frame}


%%%%%%%%%%%%%%
\section[ไฮเปอร์โบลิก]{อนุพันธ์ฟังก์ชันไฮเปอร์โบลิก}
\frame{\sectionpage}
%%%%%%%%%%%%%%

\begin{frame}
\frametitle{เส้นโค้งไฮเปอร์โบลิกคืออะไร}
พิจารณาโซ่ที่คล้องบนเสาสองข้าง ถ้าโซ่มีความยาวมากกว่าระยะห่างระหว่างเสาแล้วโซ่จะถูกแรงโน้มถ่วงดึงลงมา ทำให้เกิดเป็นรูปเส้นโค้งขึ้นมา
\addfig{0.5}{diff_hyperbolic_intro1}{เส้นโค้งไฮเปอร์โบลิกที่เกิดจากแขวนโซ่ระหว่างเสาสองต้น}{}
\end{frame}

\begin{frame}
ในทางกลับกัน การจะสร้างเส้นโค้งสะพานให้มีโครงสร้างที่แข็งแรงและกระจายน้ำหนักไปยังพื้นสองข้างได้ดีที่สุด ต้องใช้เส้นโค้งรูปแบบเดียวกันกับโซ่
\addfig{0.35}{diff_hyperbolic_intro2}{Gateway Arch (สหรัฐอเมริกา) ที่สร้างขึ้นโดยใช้ฟังก์ชันไฮเปอร์โบลิก}{}
\end{frame}

\begin{frame}
\frametitle{นิยามฟังก์ชันไฮเปอร์โบลิก}
แม้ว่าจะดูคล้ายกับเส้นโค้งของฟังก์ชันพาราโบลา $f(x) = x^2$ แต่ที่จริงแล้วเส้นโค้งนี้ขึ้นอยู่กับฟังก์ชันเลขชี้กำลัง $f(x) = e^x$ ตามนิยามดังนี้
\begin{defn}{}{}
\begin{eq}
\sinh x &=& \frac{e^x - e^{-x}}{2} \\
\cosh x &=& \frac{e^x + e^{-x}}{2} \\
\tanh x &=& \frac{\sinh x}{\cosh x} = \frac{e^x - e^{-x}}{e^x + e^{-x}} 
\end{eq}
\end{defn}
\end{frame}

\begin{frame}
\frametitle{นิยามฟังก์ชันไฮเปอร์โบลิกส่วนกลับ}
เราสามารถนิยามฟังก์ชันตรีไฮเปอร์โบลิกส่วนกลับ ได้ในทำนองเดียวกันกับฟังก์ชันตรีโกณมิติส่วนกลับได้ดังนี้
\begin{defn}{}{}
\begin{eq}
\csch x &=& \frac{1}{\sinh x} = \frac{2}{e^x - e^{-x}} \\
\sech x &=& \frac{1}{\cosh x} = \frac{2}{e^x + e^{-x}} \\
\coth x &=& \frac{1}{\tanh x} = \frac{e^x + e^{-x}}{e^x - e^{-x}}
\end{eq}
\end{defn}
\end{frame}

\begin{frame}
\addfig{0.9}{diff_hyperbolic_graph}{กราฟของฟังก์ชันไฮเปอร์โบลิกทั้ง 6 ฟังก์ชัน}{}
\end{frame}

\begin{frame}
\frametitle{อนุพันธ์ฟังก์ชันไฮเปอร์โบลิก}
เนื่องจากฟังก์ชันไฮเปอร์โบลิก สร้างขึ้นมาจากฟังก์ชันเลขชี้กำลัง เราจึงหาอนุพันธ์ได้โดยตรงดังนี้
\begin{thm}{}{}
	\begin{columns}
		\begin{column}{0.5\textwidth}
			\begin{itemize}
				\item $\dx \sinh x = \cosh x$
				\item $\dx \cosh x = \sinh x$
				\item $\dx \tanh x = (\sech x)^2$	
			\end{itemize}
		\end{column}
		\begin{column}{0.5\textwidth}
			\begin{itemize}
				\item $\dx \coth x = (\csch x)^2$
				\item $\dx \sech x = -\sech x \tanh x$
				\item $\dx \csch x = -\csch x \coth x$
			\end{itemize}
		\end{column}
	\end{columns}
\end{thm}
การพิสูจน์อนุพันธ์ของ $\sinh x$ อยู่ในบทพิสูจน์ G)
\end{frame}